\section{One company close across several accounting methods}
\label{sec:company-composition}

Separate examples can each balance while disagreeing about the cash they share.
A return settlement may draw on the same bank account that must redeem a bond;
an internal sale may increase a parent's earnings without increasing the group's
earnings. To see these interactions, consider a second synthetic company: a
parent, P, and its subsidiary, S. P owns 80\% of S at the beginning of 2025 and
controls it throughout the year. Both use USD, the same calendar and the same
accounting policies. All amounts are currency units. The opening balances are
supplied facts; the example begins after acquisition and debt issuance.

P opens with cash 500, inventory 60, plant 100, a return-recovery asset 16, an
investment in S of 80 and a deferred tax asset of 10. It owes a refund liability
of 40, matured debt of 100 and a deferred tax liability of 4. Its capital 600,
original compound-instrument equity 10 and retained earnings 12 give equity
622. S opens with cash and capital of 100. Eliminating P's investment against
its share of S's opening capital leaves non-controlling interests (NCI) of 20.
The group's opening assets are 786, liabilities 144 and equity 642.

This case has a deliberately narrow tax law: a supplied enacted rate of 25\%,
deductible wages, matching book and tax depreciation, and the return timing
specified below. Recovery of the stated deferred tax assets is supported by a
supplied finding. P controls the timing of reversal of taxable outside-basis
investment differences, with no reversal expected in the foreseeable future,
so the IAS~12.39 exception applies conditionally. These are teaching premises,
not a statement of any jurisdiction's current law. The recognition,
consolidation and cash-flow rules use the retained historical standard text
\citep[IAS~12.24,39,46--51; IFRS~10.B86--B94]{eu20231803}.

\subsection{Cash paid and expense incurred answer different questions}

The bond's original liability proceeds were 90 and its original equity component
was 10. Accrued interest of 10 has already increased the opening liability to
100. Redemption on 2 January therefore debits debt 100 and credits cash 100.
It creates no new 2025 interest expense. Under this company's historical IAS~7
operating election, the payment is financing principal 90 and operating interest
10. Interest paid must still be disclosed separately even when the associated
expense belongs to another period \citep[IAS~7.31--33]{eu20231803}.

A component adapter must consequently preserve the original debt cash basis
and the accumulated liability. Carrying amount alone cannot reconstruct the
cash-flow split. Nor can an event labelled ``accrue'' simply be moved across a
year boundary: the opening states whether its expense has already been recorded.
Here P's cash falls to 400, the debt becomes zero and original equity remains 10.

The return liability comes from a previous ten-unit sale. Each unit originally
cost 12 and can be refunded for 20. Expected returns were two units; estimated
handling cost 3 and deterioration 1 per unit gave a recovery asset
\(2(12-3-1)=16\) and refund liability \(2(20)=40\). Two units actually return.
Handling is 4 per unit, deterioration is 2 and the returned stock's supplied
net realisable value is 20. All return rights are now closed.

The stock recovered is \(\min\{2(12-2),20\}=20\). P pays refunds 40 and
handling 8, releases the refund liability 40 and recovery asset 16, and records
inventory 20. Because the actual handling is explicitly a selling expense,
the remaining entries are a cost-of-sales credit of 4 and selling expense 8.
Their net expense is 4; there is no revenue adjustment when actual refunds equal
the liability. This separates the recovery estimate from the cash actually
spent and avoids charging handling twice
\citep[IFRS~15.B21--B27; IAS~2.16]{eu20231803}.

Tax timing supplies the next link. The original sale was already taxable, and
its original inventory cost was deducted. On return, tax deducts the refund 40,
restores recovered inventory 20 to tax stock and deducts handling 8. Thus the
return reduces taxable income by 28, compared with book expense 4. The
difference is 24. Opening deferred tax asset 10 and liability 4 reverse, giving
deferred tax expense 6. A zero-tax return example would miss this interaction.

\subsection{The subsidiary's inventory contains the parent's margin}

During the first half P also pays wages 10 and records depreciation 10. It then
sells its other inventory, which cost 60, to S for 140 on credit. P recognises
revenue 140 and cost of sales 60; S records inventory and a payable of 140.
The legal entries are necessary: P and S have different owners' claims and
separate obligations. But the group has not yet sold that inventory to anyone
outside the group.

S sells two-fifths of the lot externally for 84, recording legal cost of sales
56, and pays P the full internal invoice of 140. The diagram distinguishes the
two sales. Only the second crosses the group boundary.

\begin{figure}[htbp]
\centering
\begin{tikzpicture}[>=Stealth, every node/.style={font=\small}]
\node[draw,rounded corners,minimum width=26mm,minimum height=13mm] (p) at (0,0) {Parent P};
\node[draw,rounded corners,minimum width=26mm,minimum height=13mm] (s) at (4.3,0) {Subsidiary S};
\node[draw,minimum width=28mm,minimum height=13mm] (c) at (8.6,0) {Customer};
\draw[->] (p)--node[above,align=center]{whole lot\\internal price 140} (s);
\draw[->] (s)--node[above,align=center]{two-fifths\\external price 84} (c);
\draw[dashed] (-1.7,-1.2) rectangle (6,1.3);
\node[anchor=west] at (-1.7,-1.6) {Group cost of the whole lot: 60};
\end{tikzpicture}
\caption{The internal price changes the legal books; the external sale determines group revenue.}
\end{figure}

After the external sale, S carries three-fifths of its purchase price, or 84.
The same physical stock cost the group only \(3(60)/5=36\). Its unrealised
internal profit is therefore \(84-36=48\). IFRS~10.B86(c) requires full
elimination of intragroup amounts and of profit still contained in assets; it
also directs the resulting temporary differences to IAS~12.

For clarity, let \(T\) be the internal price, \(C\) the group's original cost
and \(s\) the fraction sold externally. The remaining internal profit is
\begin{equation}
 U=(T-C)(1-s)=(140-60)(1-2/5)=48.
 \label{eq:joined-unrealised}
\end{equation}
The group removes internal revenue 140 and reduces the combined legal cost of
sales by \(T-U=92\). S's tax base remains 84, while group carrying amount is
36. The deductible difference 48 gives a deferred tax asset 12 at the supplied
25\% rate. The after-tax elimination is consequently 36. If all the stock were
sold outside the group, \(s=1\), both the remaining profit and this deferred
tax asset would be zero.

P's first-half pretax profit is \(80-4-10-10=56\). The return timing difference
reduces taxable income to 32, giving current tax 8; deferred tax expense 6
leaves P profit 42. S earns 28 before tax and 21 after current tax 7. Both
current-tax amounts are paid. Group profit is therefore
\begin{equation}
 \Pi_{H1}=42+21-48+12=27.
 \label{eq:joined-h1-profit}
\end{equation}
The downstream margin originated in P's sale. Its after-tax elimination belongs
to P's shareholders; NCI receives 20\% of S's legal profit 21, or 4.2. At June
the group has cash 511, inventory 56, plant 90 and deferred tax asset 12, totalling
669. It has no remaining external liability in this case. Equity is parent
644.8 plus NCI 24.2. These totals reconcile the same shared legal books used by
the return and debt methods.

\subsection{Buying more shares changes who owns the equity}

On 1 July P pays 20 to the outside shareholders for another 10\% of S. P now
owns 90\%, and control is unchanged. In P's separate books, its cost-policy
investment increases by 20 and cash decreases by 20. S records nothing: the
payment was between its shareholders. The group's accounting must instead
describe a change in the relative claims of P's owners and NCI.

Immediately before purchase, S's equity is 121 and NCI owns 20\%, carrying
24.2. P buys half of that interest. The NCI reduction is 12.1, not the cash
price 20. IFRS~10.23 and B96 put the difference directly in parent equity:
\begin{align}
 \Delta NCI&=-(0.90-0.80)121=-12.1,\\
 \Delta E_P&=-(20-12.1)=-7.9.
 \label{eq:joined-nci-purchase}
\end{align}
No gain, goodwill or remeasurement is introduced. IAS~7.42A--42B classifies the
group cash payment as financing, while the parent's separate purchase of an
investment is investing under its own cash-flow presentation
\citep[IFRS~10.23,B96; IAS~7.16,42A--42B; IAS~27.10]{eu20231803}.

The purchase date matters for later earnings. First-half NCI profit remains
4.2; it is not recomputed at 10\%. The new fraction applies to subsequent
profit. Recording only an ending ownership percentage would lose that history
and allocate the year's earnings incorrectly.

\subsection{The external sale releases both margin and tax}

In the second half S sells the remaining lot for 126. Its legal cost is 84,
so legal pretax profit is 42 and tax is 10.5. The group's cost is only 36.
The earlier elimination of 48 now reverses, and the associated deferred tax
asset 12 is consumed. P separately sells its recovered inventory for 40 at
cost 20 and records a further depreciation charge 10. Its second-half pretax
profit is 10 and tax is 2.5. All current taxes are paid.

The resulting group profit is
\begin{equation}
 \Pi_{H2}=(10-2.5)+(42-10.5)+48-12=75.
 \label{eq:joined-h2-profit}
\end{equation}
NCI receives 10\% of S's after-tax profit 31.5, or 3.15. The group release
again belongs to P's shareholders because it reverses P's downstream margin.
Annual NCI profit is \(4.2+3.15=7.35\), and annual parent profit is
\(102-7.35=94.65\). Closing equity can now be reconstructed without using
an asset total as a balancing plug:
\begin{align}
 E_P^{close}&=622+94.65-7.9=708.75,\\
 NCI^{close}&=20+7.35-12.1=15.25.
\end{align}
Their sum is 724, equal to closing cash 644 and plant 80. The inventory has
been sold, the return and debt obligations settled, and the temporary
differences in this case reversed.

The comparison below follows the derivation rather than replacing it. The
balance-sheet columns are closing amounts; income and cash columns are flows
over the stated interval.
\begin{center}
\begin{tabular}{lrrr}
\toprule & First half & Second half & Full year\\\midrule
External revenue &84&166&250\\
Cost of sales &20&56&76\\
Profit before tax &36&100&136\\
Current tax expense &15&13&28\\
Deferred tax expense/(benefit) &-6&12&6\\
Profit after tax &27&75&102\\
NCI share of profit &4.2&3.15&7.35\\
Operating cash flow &1&153&154\\
Financing cash flow &-90&-20&-110\\
Closing cash &511&644&644\\
Closing group equity &669&724&724\\\bottomrule
\end{tabular}
\end{center}

Cash provides an independent check on the joins. External receipts are
\(84+126+40=250\). Operating payments are refunds 40, handling 8, wages 10,
interest 10 and tax 28, totalling 96. Financing payments are debt principal 90
and the NCI purchase 20. Thus
\begin{equation}
 C^{close}=600+(250-96)-(90+20)=644.
\end{equation}
The internal receipt 140 and payment 140 are removed from their respective
gross totals. Merely observing that they net to zero would overstate both
group receipts and payments. Group adjustments never alter either entity's
legal cash ledger.

\subsection{The implementation consequence of a joined close}

Each method receives the current shared state, not its own duplicate cash
balance. The internal sale and settlement each update both legal books as one
transaction; an unmatched invoice or one-sided publication is refused. A
report must link an uninterrupted receipt sequence from its opening state to
its closing state. Consolidation recomputes the required worksheet adjustment
from legal balances and lot history at each date. It does not append the June
elimination again in December.

The companion specification, \code{CompanyComposition/1}, fixes these inputs,
events and outputs in closed records. An independent calculation reconstructs
the legal and group reports, including the adverse cases of balanced but wrong
journals, missing tax findings, changed ownership identities and partial
publication. Its two narrow migration examples preserve the original source
owners and dates; they do not rename a previous synthetic company into P.
Section~\ref{sec:company-composition-reference} gives the exchange details.

This joined case establishes the stated calculation and record design for the
supplied opening. It does not reconstruct the prior acquisition or issuance,
nor does a corporate group with no bank or insurance business specify those
sectors' combined statements. Foreign currency, wider ownership rights, actual
tax authority and the remaining accounting methods retain their own required
evidence. Those limits identify what the next case must add without obscuring
what the shared company close already demonstrates.
